% ============================================================================
% APPENDIX DC: LIT-CARD - Labor Economics & Natural Experiments
% ============================================================================
% Category: Literature Integration (LIT-)
% Language: English
% EBF Integration: Causal inference, labor markets, and policy evaluation
% ============================================================================

\chapter*{Appendix UNMAPPED_DC: David Card Research Integration}
\addcontentsline{toc}{chapter}{Appendix UNMAPPED_DC: LIT-CARD - Labor Economics \& Natural Experiments}
\label{app:card}

% ============================================================================
% HEADER BLOCK
% ============================================================================
\begin{tcolorbox}[colback=blue!5!white, colframe=blue!75!black, title=Appendix Metadata]
\textbf{Code:} DC \\
\textbf{Category:} LIT- (Literature Integration) \\
\textbf{Title:} Labor Economics \& Natural Experiments \\
\textbf{Version:} 1.0 \\
\textbf{Status:} Complete \\
\textbf{Primary Author Focus:} David Card (Nobel Prize 2021) \\
\textbf{Papers Covered:} 30 \\
\textbf{Dependencies:} DOMAIN-LABOR, METHOD-EXPERIMENT, LIT-HECKMAN (N), DOMAIN-EDUCATION \\
\textbf{EBF Connections:} Causal identification, labor market parameters, policy evaluation
\end{tcolorbox}

% ============================================================================
% CROSS-REFERENCE MAP
% ============================================================================
\begin{tcolorbox}[colback=green!5!white, colframe=green!75!black, title=Cross-Reference Map]
\textbf{Outgoing References:}
\begin{itemize}[nosep]
    \item DOMAIN-LABOR: Labor market theory and evidence
    \item DOMAIN-EDUCATION: Returns to schooling
    \item METHOD-EXPERIMENT: Natural experiment methodology
    \item LIT-HECKMAN (N): Selection and causal inference
    \item LIT-DELLAVIGNA (DV): Field experiment collaboration
    \item CONTEXT-DEMOGRAPHIC: Immigration and gender effects
\end{itemize}

\textbf{Incoming References:}
\begin{itemize}[nosep]
    \item Chapter 4 (Methodology): Design-based research
    \item Chapter 14 (Applications): Labor policy design
    \item METHOD-META: Active labor market meta-analyses
\end{itemize}
\end{tcolorbox}

% ============================================================================
% ABSTRACT
% ============================================================================
\section*{Abstract}

David Card revolutionized empirical economics through his development of \textbf{natural experiment} methodology and its application to fundamental labor market questions. His 2021 Nobel Prize recognized ``empirical contributions to labour economics.'' This appendix integrates Card's 30 key papers into the EBF framework, with particular emphasis on:
\begin{enumerate}
    \item \textbf{Minimum wage effects}: The Card-Krueger revolution
    \item \textbf{Immigration and labor markets}: The Mariel Boatlift natural experiment
    \item \textbf{Returns to education}: Instrumental variable approaches
    \item \textbf{Wage inequality}: Firm effects and labor market structure
    \item \textbf{Active labor market policies}: Meta-analytic evidence
\end{enumerate}

% ============================================================================
% QUICK REFERENCE
% ============================================================================
\begin{tcolorbox}[colback=yellow!10!white, colframe=orange!75!black, title=Quick Reference: Key Card \& Labor Market Parameters]
\begin{tabular}{llll}
\textbf{Parameter} & \textbf{Value} & \textbf{Source} & \textbf{ID} \\
\hline
Min.\ wage employment effect & $\approx 0$ & Card \& Krueger (1994) & --- \\
Immigration wage elasticity & $-0.1$ to $-0.2$ & Card (1990, 2001) & --- \\
Returns to education & 8--12\% p.a. & Card (1999, 2001) & --- \\
Firm wage premium share & 20--25\% & Card et al.\ (2013) & --- \\
Peer salary dissatisfaction & $-15\%$ & Card et al.\ (2012) & --- \\
\hline
\multicolumn{4}{l}{\textit{JEP 2026 Labor Market Power Symposium:}} \\
Monopsony wage markdown & $\mu = 0.70$--$0.80$ & Berger et al.\ (2026) & PAR-LM-006 \\
Concentrated wage share & 40\% of wages & Berger et al.\ (2026) & PAR-LM-007 \\
Welfare loss from monopsony & 7.6\% lifetime & Berger et al.\ (2026) & PAR-LM-008 \\
FTC/DOJ enforcement growth & 1 $\to$ 7 actions & Prager (2026) & PAR-LM-010 \\
Cournot markdown $\propto$ HHI & $s_i/\varepsilon$ & Prager (2026) & PAR-LM-011 \\
Noncompete prevalence & 18\% of workers & Starr (2026) & PAR-LM-012 \\
Ban wage effect & +3--4\% & Starr (2026) & PAR-LM-013 \\
Licensed workforce share & 24\% (from 5\%) & Johnson (2026) & PAR-LM-014 \\
Licensing wage premium & 15\% & Johnson (2026) & PAR-LM-015 \\
\end{tabular}
\end{tcolorbox}

% ============================================================================
% FUNDAMENTAL QUESTION
% ============================================================================
% -----------------------------------------------------------------------------
% APPENDIX SCOPE BOX
% -----------------------------------------------------------------------------

\begin{tcolorbox}[
    colback=orange!5!white,
    colframe=orange!75!black,
    title={\textbf{Appendix Scope: Ziel / In-Scope / Out-of-Scope / Constraints}},
    fonttitle=\bfseries
]

\textbf{Ziel (Objective):}
\begin{itemize}[nosep]
    \item How can we use naturally occurring variation to identify causal effects in labor markets, and what do these effects tell us about the functioning of labor market institutions?
\end{itemize}

\textbf{In-Scope:}
\begin{itemize}[nosep]
    \item Fundamental Question
    \item Minimum Wage and Employment
    \item Immigration and Labor Markets
    \item Returns to Education
    \item Wage Inequality and Firm Effects
\end{itemize}

\textbf{Out-of-Scope (delegiert an andere Appendices):}
\begin{itemize}[nosep]
    \item Glossary $\rightarrow$ Appendix G
\end{itemize}

\textbf{Constraints (Anwendungsgrenzen):}
\begin{itemize}[nosep]
    \item [TODO: Define when this appendix does NOT apply]
    \item [TODO: Define required prerequisites]
    \item [TODO: Define known limitations]
\end{itemize}

\textbf{Lieferobjekte (Deliverables):}
\begin{itemize}[nosep]
    \item 1 Table(s)
    \item 6 Equation(s)
    \item Worked Example(s)
\end{itemize}

\end{tcolorbox}


\section{Fundamental Question}

\begin{tcolorbox}[colback=red!5!white, colframe=red!75!black, title=Central Question]
\textbf{How can we use naturally occurring variation to identify causal effects in labor markets, and what do these effects tell us about the functioning of labor market institutions?}
\end{tcolorbox}

Card's answer: By exploiting \textbf{natural experiments}---policy changes, geographic variation, and institutional discontinuities---we can identify causal effects that challenge conventional economic wisdom and inform evidence-based policy.

% ============================================================================
% THEORETICAL FOUNDATIONS
% ============================================================================
\section{Theoretical Foundations: Design-Based Research}

\subsection{The Natural Experiment Revolution}

Card's methodological contribution (Nobel Lecture 2022) emphasizes:

\begin{definition}[Design-Based Research]
An approach to causal inference that:
\begin{enumerate}
    \item Exploits naturally occurring variation (policy changes, geographic borders, cohort differences)
    \item Focuses on \textbf{identification} before estimation
    \item Uses transparent assumptions that can be tested
    \item Prioritizes internal validity with clear scope conditions
\end{enumerate}
\end{definition}

\begin{axiom}[UNMAPPED_DC-1: Natural Experiment Validity]
\label{ax:dc1}
A natural experiment provides valid causal estimates when:
\begin{equation}
E[Y_0 | D=1] = E[Y_0 | D=0]
\end{equation}
i.e., the treatment assignment is ``as good as random'' conditional on observables.
\textbf{EBF Connection:} Foundation for METHOD-EXPERIMENT causal identification.
\end{axiom}

% ============================================================================
% MINIMUM WAGE
% ============================================================================
\section{Minimum Wage and Employment}

\subsection{The Card-Krueger Revolution}

Card \& Krueger (1994) challenged the textbook prediction that minimum wages reduce employment:

\begin{tcolorbox}[colback=blue!5!white, colframe=blue!75!black]
\textbf{The New Jersey Natural Experiment:}
\begin{itemize}
    \item New Jersey raised minimum wage from \$4.25 to \$5.05 in 1992
    \item Pennsylvania (control) kept minimum wage at \$4.25
    \item Fast-food employment \textbf{increased} in NJ relative to PA
    \item Difference-in-differences estimate: +2.76 FTE employees per restaurant
\end{itemize}
\end{tcolorbox}

\begin{axiom}[UNMAPPED_DC-2: Monopsony in Labor Markets]
\label{ax:dc2}
The near-zero employment effect of minimum wages suggests labor market monopsony:
\begin{equation}
\frac{\partial L}{\partial w_{min}} \approx 0 \quad \text{when } w_{min} < w^{monopsony}
\end{equation}
Firms have wage-setting power, so moderate minimum wage increases transfer rents without reducing employment.
\textbf{EBF Connection:} Challenges simple supply-demand models in DOMAIN-LABOR.
\end{axiom}

\subsection{Policy Implications}

Card \& Krueger (1995) ``Myth and Measurement'' synthesizes the evidence:
\begin{itemize}
    \item Small minimum wage increases have minimal employment effects
    \item Larger increases may have negative effects at the margin
    \item Effects vary by labor market tightness and firm characteristics
    \item Distributional effects favor low-wage workers
\end{itemize}

% ============================================================================
% IMMIGRATION
% ============================================================================
\section{Immigration and Labor Markets}

\subsection{The Mariel Boatlift}

Card (1990) exploited the sudden influx of 125,000 Cuban immigrants to Miami:

\begin{tcolorbox}[colback=green!5!white, colframe=green!75!black]
\textbf{Mariel Boatlift Natural Experiment:}
\begin{itemize}
    \item 7\% increase in Miami labor force in 1980
    \item Comparison cities: Atlanta, Houston, Los Angeles, Tampa
    \item \textbf{No significant effect} on wages or unemployment of natives
    \item Even for low-skilled workers most similar to Mariel immigrants
\end{itemize}
\end{tcolorbox}

\begin{axiom}[UNMAPPED_DC-3: Labor Market Absorption]
\label{ax:dc3}
Labor markets can absorb large immigrant inflows through:
\begin{enumerate}
    \item Demand adjustments (industry composition shifts)
    \item Native out-migration to other regions
    \item Capital deepening in response to labor supply
\end{enumerate}
\begin{equation}
\frac{\partial w_{native}}{\partial L_{immigrant}} \approx -0.1 \text{ to } -0.2
\end{equation}
\textbf{EBF Connection:} Informs CONTEXT-DEMOGRAPHIC parameters for policy.
\end{axiom}

\subsection{Places vs. People: Reconciling the Literature}

\citet{dustmann2025immigration} provide a unified framework that reconciles the seemingly contradictory findings between regional (Card-style) and worker-level (Borjas-style) approaches:

\begin{tcolorbox}[colback=blue!5!white, colframe=blue!75!black]
\textbf{Key Insight: Both Approaches Are Valid But Answer Different Questions}
\begin{itemize}
    \item \textbf{Regional approaches} (Card, 1990): Measure effects on \textit{places}---including displaced natives who leave
    \item \textbf{Worker-level approaches} (Borjas, 2003): Measure effects on \textit{people}---natives who remain in competition
\end{itemize}
\end{tcolorbox}

\textbf{Decomposition of Regional Employment Effects:}
\begin{equation}
\frac{dL^N_r}{dL^I_r} = \underbrace{\text{Displacement}}_{\approx 0.14} + \underbrace{\text{Crowding-out}}_{\approx 0.02} + \underbrace{\text{Relocation}}_{\approx 0.71} = -0.87
\end{equation}

\textbf{Key Parameters from German Data:}
\begin{itemize}
    \item Inverse labor demand elasticity: $\phi^{-1} = -1.95$
    \item Place-level labor supply elasticity: $\eta^P = 4.64$
    \item Worker-level labor supply elasticity: $\eta^E = 3.68$
    \item Pure wage effect: $-0.188$
\end{itemize}

\textbf{Policy Implication:} Neither approach is ``correct''---both provide useful estimates for different policy questions. Regional effects inform area-level policies; worker-level effects inform individual welfare assessments.

% ============================================================================
% RETURNS TO EDUCATION
% ============================================================================
\section{Returns to Education}

\subsection{The Identification Challenge}

Card (1999, 2001) addresses the endogeneity of schooling:

\begin{itemize}
    \item \textbf{Ability bias}: More able individuals get more education
    \item \textbf{Measurement error}: Self-reported education is noisy
    \item \textbf{Selection}: Who chooses to continue schooling?
\end{itemize}

\subsection{Instrumental Variable Solutions}

\begin{axiom}[UNMAPPED_DC-4: IV Estimates of Returns]
\label{ax:dc4}
Using instruments (college proximity, compulsory schooling laws, twins), the causal return to education is:
\begin{equation}
\frac{\partial \ln(wage)}{\partial years\_education} \approx 0.08-0.12
\end{equation}
IV estimates are often \textbf{higher} than OLS, suggesting ability bias is offset by measurement error attenuation.
\textbf{EBF Connection:} Provides calibration for DOMAIN-EDUCATION human capital models.
\end{axiom}

% ============================================================================
% WAGE INEQUALITY
% ============================================================================
\section{Wage Inequality and Firm Effects}

\subsection{The Role of Firms}

Card, Heining \& Kline (2013) decompose wage inequality using matched employer-employee data:

\begin{equation}
\ln(w_{ijt}) = \alpha_i + \psi_j + X'_{it}\beta + \varepsilon_{ijt}
\end{equation}

where $\alpha_i$ is worker fixed effect, $\psi_j$ is firm fixed effect.

\begin{axiom}[UNMAPPED_DC-5: Firm Wage Premiums]
\label{ax:dc5}
Firm effects explain 20-25\% of wage variance:
\begin{itemize}
    \item High-wage firms pay \textbf{all} workers more (not just high-skill)
    \item Sorting of high-skill workers to high-paying firms increased over time
    \item Rising inequality = rising firm dispersion + increased sorting
\end{itemize}
\textbf{EBF Connection:} Extends DOMAIN-LABOR with firm-level heterogeneity.
\end{axiom}

\subsection{Peer Effects and Fairness}

Card, Mas, Moretti \& Saez (2012) show that relative pay matters:

\begin{tcolorbox}[colback=yellow!10!white, colframe=yellow!75!black]
\textbf{Peer Salary Experiment:}
\begin{itemize}
    \item UC employees given access to colleagues' salaries
    \item Those earning \textbf{below} median: job satisfaction drops 15\%
    \item Those earning \textbf{above} median: no change in satisfaction
    \item Asymmetric response consistent with loss aversion
\end{itemize}
\end{tcolorbox}

\begin{axiom}[UNMAPPED_DC-6: Social Comparison in Wages]
\label{ax:dc6}
Utility from wages includes a social comparison component:
\begin{equation}
U(w_i) = u(w_i) + \lambda \cdot \mathbb{1}_{w_i < w_{reference}} \cdot (w_{reference} - w_i)
\end{equation}
where $\lambda > 0$ captures loss aversion to being below reference.
\textbf{EBF Connection:} Links to CORE-WHAT fairness dimension and loss aversion.
\end{axiom}

% ============================================================================
% ACTIVE LABOR MARKET POLICY
% ============================================================================
\section{Active Labor Market Policy}

\subsection{Meta-Analytic Evidence}

Card, Kluve \& Weber (2010, 2018) synthesize evidence from 200+ program evaluations:

\begin{tcolorbox}[colback=blue!5!white, colframe=blue!75!black, title=What Works in Labor Market Policy?]
\begin{tabular}{lcc}
\textbf{Program Type} & \textbf{Short-run} & \textbf{Long-run} \\
\hline
Job search assistance & ++ & + \\
Training programs & 0/+ & ++ \\
Private sector subsidies & + & + \\
Public employment & 0 & 0 \\
\end{tabular}
\begin{itemize}
    \item Training programs: negative short-run (lock-in), positive long-run
    \item Public employment programs: consistently ineffective
    \item Effects larger for women, long-term unemployed
\end{itemize}
\end{tcolorbox}

% ============================================================================
% EBF INTEGRATION
% ============================================================================
\section{EBF Integration}

\subsection{10C Mapping}

\begin{table}[h]
\centering
\begin{tabular}{lll}
\textbf{10C Element} & \textbf{Card Contribution} & \textbf{Parameter} \\
\hline
CORE-WHO (AAA) & Worker/firm heterogeneity & Firm effects $\psi_j$ \\
CORE-WHAT (C) & Fairness in wages & Social comparison $\lambda$ \\
CORE-HOW (B) & Labor market complementarities & Skill-firm matching, $\gamma_{granularity}$ \\
CORE-WHEN (V) & Policy timing \& institutional context & Short vs. long-run, $\Psi_I$ \\
CORE-WHERE (BBB) & Labor market parameters & Returns, elasticities, PAR-LM-006--009 \\
CORE-AWARE (AU) & Wage transparency effects & Peer salary knowledge \\
CORE-READY (AV) & Job search intensity & ALMP effectiveness \\
CORE-STAGE (AW) & Career progression & Education-wage path \\
CORE-HIERARCHY (HI) & Firm-level $\to$ market power & Granularity channel \\
\end{tabular}
\caption{Card contributions mapped to 10C framework}
\end{table}

\subsection{Key Parameter Values for BBB}

\begin{tcolorbox}[colback=green!5!white, colframe=green!75!black, title=Card Parameters for BBB Repository]
\begin{itemize}
    \item Returns to education: 8-12\% per year (IV estimates)
    \item Minimum wage employment elasticity: $\approx 0$ (for moderate increases)
    \item Immigration wage elasticity: -0.1 to -0.2
    \item Firm wage premium variance share: 20-25\%
    \item Peer salary dissatisfaction: -15\% if below median
    \item Training program long-run effect: +5-10\% earnings
    \item \textbf{Monopsony wage markdown: $\mu = 0.70$--$0.80$ (PAR-LM-006)}
    \item \textbf{Labor market HHI threshold: 2,500; 40\% of wages in concentrated markets (PAR-LM-007)}
    \item \textbf{Welfare loss from monopsony: 7.6\% lifetime consumption (PAR-LM-008)}
    \item \textbf{Noncompete wage effect: +4\% from elimination (PAR-LM-009)}
\end{itemize}
\end{tcolorbox}

% ============================================================================
% WORKED EXAMPLE
% ============================================================================
\section{Worked Example: Minimum Wage Policy Analysis}

Consider a city contemplating a minimum wage increase from \$12 to \$15:

\textbf{Card Framework Analysis:}

\begin{enumerate}
    \item \textbf{Identify comparison group}: Neighboring city without increase
    \item \textbf{Estimate baseline}: Current employment in affected industries
    \item \textbf{Apply Card-Krueger elasticity}: $\epsilon \approx 0$ for moderate increase
    \item \textbf{Predict distributional effects}:
    \begin{itemize}
        \item Workers earning \$12-15: wage increase to \$15
        \item Workers earning >\$15: potential spillover (+2-5\%)
        \item Employment: minimal change predicted
    \end{itemize}
    \item \textbf{Monitor with difference-in-differences}: Compare trends post-policy
\end{enumerate}

\textbf{Caveat}: Large increases (>30\%) may show negative employment effects.

% ============================================================================
% SUMMARY
% ============================================================================


%%%%%%%%%%%%%%%%%%%%%%%%%%%%%%%%%%%%%%%%%%%%%%%%%%%%%%%%%%%%%%%%%%%%%%%%%%%%%%%
\section{Key Results}
\label{sec:dc-results}
%%%%%%%%%%%%%%%%%%%%%%%%%%%%%%%%%%%%%%%%%%%%%%%%%%%%%%%%%%%%%%%%%%%%%%%%%%%%%%%

The analysis yields the following key results:

\begin{enumerate}
    \item \textbf{Result 1:} Primary finding with quantitative support
    \item \textbf{Result 2:} Secondary finding with implications
    \item \textbf{Result 3:} Practical application or boundary condition
\end{enumerate}


% === AUTO-GENERATED ENTRIES (2026-02-22) ===

%%%%%%%%%%%%%%%%%%%%%%%%%%%%%%%%%%%%%%%%%%%%%%%%%%%%%%%%%%%%%%%%%%%%%%%%%%%%%%%
\subsubsection{2013: Nine Facts about Top Journals in Economics}
%%%%%%%%%%%%%%%%%%%%%%%%%%%%%%%%%%%%%%%%%%%%%%%%%%%%%%%%%%%%%%%%%%%%%%%%%%%%%%%

\paragraph{Citation.}
Card, David and DellaVigna, Stefano (2013). ``Nine Facts about Top Journals in Economics.'' \textit{Journal of Economic Literature}, 51, 144--161.

\paragraph{10C Integration.}
\begin{itemize}[nosep]
  \item \textbf{Domain}: [To be specified]
  \item \textbf{use\_for}: LIT-CARD
\end{itemize}


%%%%%%%%%%%%%%%%%%%%%%%%%%%%%%%%%%%%%%%%%%%%%%%%%%%%%%%%%%%%%%%%%%%%%%%%%%%%%%%
\subsubsection{2020: What Do Editors Maximize? Evidence from Four Economics Journ...}
%%%%%%%%%%%%%%%%%%%%%%%%%%%%%%%%%%%%%%%%%%%%%%%%%%%%%%%%%%%%%%%%%%%%%%%%%%%%%%%

\paragraph{Citation.}
Card, David and DellaVigna, Stefano (2020). ``What Do Editors Maximize? Evidence from Four Economics Journals.'' \textit{Review of Economics and Statistics}, 102, 195--217.

\paragraph{10C Integration.}
\begin{itemize}[nosep]
  \item \textbf{Domain}: [To be specified]
  \item \textbf{use\_for}: LIT-CARD
\end{itemize}


%%%%%%%%%%%%%%%%%%%%%%%%%%%%%%%%%%%%%%%%%%%%%%%%%%%%%%%%%%%%%%%%%%%%%%%%%%%%%%%
\subsubsection{1999: The Return to Schooling in Canada}
%%%%%%%%%%%%%%%%%%%%%%%%%%%%%%%%%%%%%%%%%%%%%%%%%%%%%%%%%%%%%%%%%%%%%%%%%%%%%%%

\paragraph{Citation.}
Card, David (1999). ``The Return to Schooling in Canada.'' \textit{Canadian Journal of Economics}, 32.

\paragraph{10C Integration.}
\begin{itemize}[nosep]
  \item \textbf{Domain}: [To be specified]
  \item \textbf{use\_for}: LIT-CARD
\end{itemize}


%%%%%%%%%%%%%%%%%%%%%%%%%%%%%%%%%%%%%%%%%%%%%%%%%%%%%%%%%%%%%%%%%%%%%%%%%%%%%%%
\subsubsection{1992: Does School Quality Matter? Returns to Education and the Cha...}
%%%%%%%%%%%%%%%%%%%%%%%%%%%%%%%%%%%%%%%%%%%%%%%%%%%%%%%%%%%%%%%%%%%%%%%%%%%%%%%

\paragraph{Citation.}
Card, David and Krueger, Alan (1992). ``Does School Quality Matter? Returns to Education and the Characteristics of Public Schools in the United States.'' \textit{Journal of Political Economy}, 100.

\paragraph{10C Integration.}
\begin{itemize}[nosep]
  \item \textbf{Domain}: [To be specified]
  \item \textbf{use\_for}: LIT-CARD
\end{itemize}


%%%%%%%%%%%%%%%%%%%%%%%%%%%%%%%%%%%%%%%%%%%%%%%%%%%%%%%%%%%%%%%%%%%%%%%%%%%%%%%
\subsubsection{1995: Using Geographic Variation in College Proximity to Estimate ...}
%%%%%%%%%%%%%%%%%%%%%%%%%%%%%%%%%%%%%%%%%%%%%%%%%%%%%%%%%%%%%%%%%%%%%%%%%%%%%%%

\paragraph{Citation.}
Card, David (1995). ``Using Geographic Variation in College Proximity to Estimate the Return to Schooling.'' \textit{Aspects of Labour Market Behaviour}.

\paragraph{10C Integration.}
\begin{itemize}[nosep]
  \item \textbf{Domain}: [To be specified]
  \item \textbf{use\_for}: LIT-CARD
\end{itemize}


%%%%%%%%%%%%%%%%%%%%%%%%%%%%%%%%%%%%%%%%%%%%%%%%%%%%%%%%%%%%%%%%%%%%%%%%%%%%%%%
\subsubsection{1996: Wage Dispersion, Returns to Skill, and Black-White Wage Diff...}
%%%%%%%%%%%%%%%%%%%%%%%%%%%%%%%%%%%%%%%%%%%%%%%%%%%%%%%%%%%%%%%%%%%%%%%%%%%%%%%

\paragraph{Citation.}
Card, David and Lemieux, Thomas (1996). ``Wage Dispersion, Returns to Skill, and Black-White Wage Differentials.'' \textit{Journal of Econometrics}, 74.

\paragraph{10C Integration.}
\begin{itemize}[nosep]
  \item \textbf{Domain}: [To be specified]
  \item \textbf{use\_for}: LIT-CARD
\end{itemize}


%%%%%%%%%%%%%%%%%%%%%%%%%%%%%%%%%%%%%%%%%%%%%%%%%%%%%%%%%%%%%%%%%%%%%%%%%%%%%%%
\subsubsection{1996: The Effect of Unions on the Structure of Wages: A Longitudin...}
%%%%%%%%%%%%%%%%%%%%%%%%%%%%%%%%%%%%%%%%%%%%%%%%%%%%%%%%%%%%%%%%%%%%%%%%%%%%%%%

\paragraph{Citation.}
Card, David (1996). ``The Effect of Unions on the Structure of Wages: A Longitudinal Analysis.'' \textit{Econometrica}, 64.

\paragraph{10C Integration.}
\begin{itemize}[nosep]
  \item \textbf{Domain}: [To be specified]
  \item \textbf{use\_for}: LIT-CARD
\end{itemize}


%%%%%%%%%%%%%%%%%%%%%%%%%%%%%%%%%%%%%%%%%%%%%%%%%%%%%%%%%%%%%%%%%%%%%%%%%%%%%%%
\subsubsection{2001: Immigrant Inflows, Native Outflows, and the Local Labor Mark...}
%%%%%%%%%%%%%%%%%%%%%%%%%%%%%%%%%%%%%%%%%%%%%%%%%%%%%%%%%%%%%%%%%%%%%%%%%%%%%%%

\paragraph{Citation.}
Card, David (2001). ``Immigrant Inflows, Native Outflows, and the Local Labor Market Impacts of Higher Immigration.'' \textit{Journal of Labor Economics}, 19.

\paragraph{10C Integration.}
\begin{itemize}[nosep]
  \item \textbf{Domain}: [To be specified]
  \item \textbf{use\_for}: LIT-CARD
\end{itemize}


%%%%%%%%%%%%%%%%%%%%%%%%%%%%%%%%%%%%%%%%%%%%%%%%%%%%%%%%%%%%%%%%%%%%%%%%%%%%%%%
\subsubsection{2002: Skill-Biased Technological Change and Rising Wage Inequality...}
%%%%%%%%%%%%%%%%%%%%%%%%%%%%%%%%%%%%%%%%%%%%%%%%%%%%%%%%%%%%%%%%%%%%%%%%%%%%%%%

\paragraph{Citation.}
Card, David and DiNardo, John (2002). ``Skill-Biased Technological Change and Rising Wage Inequality: Some Problems and Puzzles.'' \textit{Journal of Labor Economics}, 20.

\paragraph{10C Integration.}
\begin{itemize}[nosep]
  \item \textbf{Domain}: [To be specified]
  \item \textbf{use\_for}: LIT-CARD
\end{itemize}


%%%%%%%%%%%%%%%%%%%%%%%%%%%%%%%%%%%%%%%%%%%%%%%%%%%%%%%%%%%%%%%%%%%%%%%%%%%%%%%
\subsubsection{2004: Unions and Wage Inequality}
%%%%%%%%%%%%%%%%%%%%%%%%%%%%%%%%%%%%%%%%%%%%%%%%%%%%%%%%%%%%%%%%%%%%%%%%%%%%%%%

\paragraph{Citation.}
Card, David and Lemieux, Thomas and Riddell, W. Craig (2004). ``Unions and Wage Inequality.'' \textit{Journal of Labor Research}, 25.

\paragraph{10C Integration.}
\begin{itemize}[nosep]
  \item \textbf{Domain}: [To be specified]
  \item \textbf{use\_for}: LIT-CARD
\end{itemize}


%%%%%%%%%%%%%%%%%%%%%%%%%%%%%%%%%%%%%%%%%%%%%%%%%%%%%%%%%%%%%%%%%%%%%%%%%%%%%%%
\subsubsection{2001: Can Falling Supply Explain the Rising Return to College for ...}
%%%%%%%%%%%%%%%%%%%%%%%%%%%%%%%%%%%%%%%%%%%%%%%%%%%%%%%%%%%%%%%%%%%%%%%%%%%%%%%

\paragraph{Citation.}
Card, David and Lemieux, Thomas (2001). ``Can Falling Supply Explain the Rising Return to College for Younger Men?.'' \textit{Quarterly Journal of Economics}, 116.

\paragraph{10C Integration.}
\begin{itemize}[nosep]
  \item \textbf{Domain}: [To be specified]
  \item \textbf{use\_for}: LIT-CARD
\end{itemize}


%%%%%%%%%%%%%%%%%%%%%%%%%%%%%%%%%%%%%%%%%%%%%%%%%%%%%%%%%%%%%%%%%%%%%%%%%%%%%%%
\subsubsection{2016: Bargaining, Sorting, and the Gender Wage Gap: Quantifying th...}
%%%%%%%%%%%%%%%%%%%%%%%%%%%%%%%%%%%%%%%%%%%%%%%%%%%%%%%%%%%%%%%%%%%%%%%%%%%%%%%

\paragraph{Citation.}
Card, David and Cardoso, Ana Rute and Kline, Patrick (2016). ``Bargaining, Sorting, and the Gender Wage Gap: Quantifying the Impact of Firms on the Relative Pay of Women.'' \textit{Quarterly Journal of Economics}, 131.

\paragraph{10C Integration.}
\begin{itemize}[nosep]
  \item \textbf{Domain}: [To be specified]
  \item \textbf{use\_for}: LIT-CARD
\end{itemize}


%%%%%%%%%%%%%%%%%%%%%%%%%%%%%%%%%%%%%%%%%%%%%%%%%%%%%%%%%%%%%%%%%%%%%%%%%%%%%%%
\subsubsection{2008: Tipping and the Dynamics of Segregation}
%%%%%%%%%%%%%%%%%%%%%%%%%%%%%%%%%%%%%%%%%%%%%%%%%%%%%%%%%%%%%%%%%%%%%%%%%%%%%%%

\paragraph{Citation.}
Card, David and Mas, Alexandre and Rothstein, Jesse (2008). ``Tipping and the Dynamics of Segregation.'' \textit{Quarterly Journal of Economics}, 123.

\paragraph{10C Integration.}
\begin{itemize}[nosep]
  \item \textbf{Domain}: [To be specified]
  \item \textbf{use\_for}: LIT-CARD
\end{itemize}


%%%%%%%%%%%%%%%%%%%%%%%%%%%%%%%%%%%%%%%%%%%%%%%%%%%%%%%%%%%%%%%%%%%%%%%%%%%%%%%
\subsubsection{2009: Does Medicare Save Lives?}
%%%%%%%%%%%%%%%%%%%%%%%%%%%%%%%%%%%%%%%%%%%%%%%%%%%%%%%%%%%%%%%%%%%%%%%%%%%%%%%

\paragraph{Citation.}
Card, David and Dobkin, Carlos and Maestas, Nicole (2009). ``Does Medicare Save Lives?.'' \textit{Quarterly Journal of Economics}, 124.

\paragraph{10C Integration.}
\begin{itemize}[nosep]
  \item \textbf{Domain}: [To be specified]
  \item \textbf{use\_for}: LIT-CARD
\end{itemize}


%%%%%%%%%%%%%%%%%%%%%%%%%%%%%%%%%%%%%%%%%%%%%%%%%%%%%%%%%%%%%%%%%%%%%%%%%%%%%%%
\subsubsection{2009: Immigration and Inequality}
%%%%%%%%%%%%%%%%%%%%%%%%%%%%%%%%%%%%%%%%%%%%%%%%%%%%%%%%%%%%%%%%%%%%%%%%%%%%%%%

\paragraph{Citation.}
Card, David (2009). ``Immigration and Inequality.'' \textit{American Economic Review}, 99.

\paragraph{10C Integration.}
\begin{itemize}[nosep]
  \item \textbf{Domain}: [To be specified]
  \item \textbf{use\_for}: LIT-CARD
\end{itemize}


%%%%%%%%%%%%%%%%%%%%%%%%%%%%%%%%%%%%%%%%%%%%%%%%%%%%%%%%%%%%%%%%%%%%%%%%%%%%%%%
\subsubsection{2012: Immigration, Wages, and Compositional Amenities}
%%%%%%%%%%%%%%%%%%%%%%%%%%%%%%%%%%%%%%%%%%%%%%%%%%%%%%%%%%%%%%%%%%%%%%%%%%%%%%%

\paragraph{Citation.}
Card, David and Dustmann, Christian and Preston, Ian (2012). ``Immigration, Wages, and Compositional Amenities.'' \textit{Journal of the European Economic Association}, 10.

\paragraph{10C Integration.}
\begin{itemize}[nosep]
  \item \textbf{Domain}: [To be specified]
  \item \textbf{use\_for}: LIT-CARD, LIT-DUSTMANN
\end{itemize}


%%%%%%%%%%%%%%%%%%%%%%%%%%%%%%%%%%%%%%%%%%%%%%%%%%%%%%%%%%%%%%%%%%%%%%%%%%%%%%%
\subsubsection{2018: Firms and Labor Market Inequality: Evidence and Some Theory}
%%%%%%%%%%%%%%%%%%%%%%%%%%%%%%%%%%%%%%%%%%%%%%%%%%%%%%%%%%%%%%%%%%%%%%%%%%%%%%%

\paragraph{Citation.}
Card, David and Cardoso, Ana Rute and Heining, Jörg and Kline, Patrick (2018). ``Firms and Labor Market Inequality: Evidence and Some Theory.'' \textit{Journal of Labor Economics}, 36.

\paragraph{10C Integration.}
\begin{itemize}[nosep]
  \item \textbf{Domain}: [To be specified]
  \item \textbf{use\_for}: LIT-CARD
\end{itemize}


%%%%%%%%%%%%%%%%%%%%%%%%%%%%%%%%%%%%%%%%%%%%%%%%%%%%%%%%%%%%%%%%%%%%%%%%%%%%%%%
\subsubsection{2020: Are Referees and Editors in Economics Gender Neutral?}
%%%%%%%%%%%%%%%%%%%%%%%%%%%%%%%%%%%%%%%%%%%%%%%%%%%%%%%%%%%%%%%%%%%%%%%%%%%%%%%

\paragraph{Citation.}
Card, David and DellaVigna, Stefano and Funk, Patricia and Iriberri, Nagore (2020). ``Are Referees and Editors in Economics Gender Neutral?.'' \textit{Quarterly Journal of Economics}, 135.

\paragraph{10C Integration.}
\begin{itemize}[nosep]
  \item \textbf{Domain}: [To be specified]
  \item \textbf{use\_for}: LIT-CARD
\end{itemize}


%%%%%%%%%%%%%%%%%%%%%%%%%%%%%%%%%%%%%%%%%%%%%%%%%%%%%%%%%%%%%%%%%%%%%%%%%%%%%%%
\subsubsection{2007: Racial Segregation and the Black-White Test Score Gap}
%%%%%%%%%%%%%%%%%%%%%%%%%%%%%%%%%%%%%%%%%%%%%%%%%%%%%%%%%%%%%%%%%%%%%%%%%%%%%%%

\paragraph{Citation.}
Card, David and Rothstein, Jesse (2007). ``Racial Segregation and the Black-White Test Score Gap.'' \textit{Journal of Public Economics}, 91.

\paragraph{10C Integration.}
\begin{itemize}[nosep]
  \item \textbf{Domain}: [To be specified]
  \item \textbf{use\_for}: LIT-CARD
\end{itemize}


%%%%%%%%%%%%%%%%%%%%%%%%%%%%%%%%%%%%%%%%%%%%%%%%%%%%%%%%%%%%%%%%%%%%%%%%%%%%%%%
\subsubsection{2022: Gender Differences in Peer Review}
%%%%%%%%%%%%%%%%%%%%%%%%%%%%%%%%%%%%%%%%%%%%%%%%%%%%%%%%%%%%%%%%%%%%%%%%%%%%%%%

\paragraph{Citation.}
Card, David and DellaVigna, Stefano and Funk, Patricia and Iriberri, Nagore (2022). ``Gender Differences in Peer Review.'' \textit{Working Paper}.

\paragraph{10C Integration.}
\begin{itemize}[nosep]
  \item \textbf{Domain}: [To be specified]
  \item \textbf{use\_for}: LIT-CARD
\end{itemize}


% ============================================================================
% LABOR MARKET POWER & MONOPSONY (Level 5 Integration: berger2026labor)
% ============================================================================
\section{Labor Market Power: The Granular Monopsony Framework}
\label{sec:dc-monopsony}

\subsection{Berger, Herkenhoff \& Mongey (2026): From Micro to Macro}

\textbf{Citation.} Berger, David and Herkenhoff, Kyle and Mongey, Simon (2026).
``Labor Market Power: From Micro Evidence to Macro Consequences.''
\textit{Journal of Economic Perspectives}, 40(1), 93--114.
DOI: 10.1257/jep.20251456. \cite{berger2026labor}

\textbf{Full text:} \texttt{data/paper-texts/PAP-berger2026labor.md} (L3, 10,224 words).

\begin{tcolorbox}[colback=blue!5!white, colframe=blue!75!black, title=Three Mechanisms of Labor Market Power]
The paper identifies three channels through which employer concentration distorts labor markets:
\begin{enumerate}
    \item \textbf{Granularity:} Large firms' wage decisions move market-level wages because each firm has non-negligible market share $s_i$. A 1\% wage reduction at a firm with share $s_i$ lowers the market wage by approximately $s_i\%$.
    \item \textbf{Strategic Interaction:} Firms internalize that their wage offers affect competitors' wages, leading to oligopsonistic coordination that further depresses wages below the competitive level.
    \item \textbf{Misallocation:} Workers are allocated to the ``wrong'' firms (too many at high-markdown firms, too few at productive firms), causing output and welfare losses beyond the direct wage markdown.
\end{enumerate}
\end{tcolorbox}

\begin{axiom}[UNMAPPED\_DC-7: Granular Monopsony Wage Markdown]
\label{ax:dc7}
In concentrated labor markets, firms pay workers a fraction of their marginal revenue product:
\begin{equation}
w_i = \mu_i \cdot MRP_i, \quad \mu_i \in [0.65, 0.80]
\end{equation}
where $\mu_i$ is the wage markdown (PAR-LM-006). The median US manufacturing firm pays 70--80\% of MRP, implying a 20--30\% markdown. In the most concentrated markets (HHI $\geq$ 2500), markdowns are substantially larger.

\textbf{EBF Connection:} Extends UNMAPPED\_DC-2 (Card-Krueger monopsony) with structural estimation.
Links to CORE-HOW: the complementarity between firm size ($s_i$) and wage-setting power is
nonlinear---large firms exert disproportionate market influence (MS-LM-006).
Parameter: PAR-LM-006 ($w/MRP = 0.70$--$0.80$).
\end{axiom}

\subsection{Concentration and the HHI Threshold}

Berger et al.\ define labor market concentration using the Herfindahl-Hirschman Index (HHI):
\begin{equation}
HHI_L = \sum_{j=1}^{N} s_j^2, \quad s_j = \frac{L_j}{\sum_k L_k}
\end{equation}
where $s_j$ is firm $j$'s employment share.

\begin{tcolorbox}[colback=yellow!10!white, colframe=yellow!75!black, title=Key Finding: Concentration Is Pervasive]
\begin{itemize}
    \item 40\% of US wage payments occur in markets with $HHI \geq 2500$ (PAR-LM-007)
    \item Average employer has a coworker-weighted $HHI$ of 3,200
    \item Even seemingly ``competitive'' metro areas contain highly concentrated occupation-specific markets
    \item Cross-country evidence (Denmark, Austria, Norway) shows similar concentration patterns
\end{itemize}
\end{tcolorbox}

\begin{axiom}[UNMAPPED\_DC-8: Macro Amplification of Monopsony]
\label{ax:dc8}
The aggregate welfare loss from labor market power substantially exceeds atomistic-firm benchmarks:
\begin{equation}
\Delta W_{aggregate} = 7.6\% \text{ of lifetime consumption} \approx 7 \times \Delta W_{atomistic}
\end{equation}
This amplification arises because granularity and strategic interaction create losses beyond those predicted by models that treat firms as infinitesimally small (MS-LM-006).

Decomposition of the output gap (20.9\% of efficient output):
\begin{itemize}
    \item 60\% from \textbf{misallocation} (workers at wrong firms)
    \item 40\% from \textbf{underemployment} (too few workers hired overall)
\end{itemize}

\textbf{EBF Connection:} The complementarity between granularity and strategic interaction ($\gamma > 0$)
means these mechanisms reinforce each other---eliminating only one channel removes far less than half the welfare loss.
Parameters: PAR-LM-008 ($W_{loss} = 7.6\%$).
\end{axiom}

\subsection{Policy Implications: Noncompete Clauses and Antitrust}

\begin{axiom}[UNMAPPED\_DC-9: Policy Interventions in Concentrated Labor Markets]
\label{ax:dc9}
Interventions that reduce labor market frictions or employer power generate wage gains:
\begin{itemize}
    \item Eliminating noncompete clauses: average wage increase $\approx$ 4\% (PAR-LM-009)
    \item Eliminating monopsony power in most concentrated markets: wage gains $>$ 30\%
    \item Minimum wage increases: particularly effective in concentrated markets (extends UNMAPPED\_DC-2)
\end{itemize}
\textbf{EBF Connection:} Links CORE-WHEN ($\Psi_I$ = institutional context) to intervention effectiveness.
In concentrated markets ($\Psi_I$ = ``few dominant employers''), labor market regulations have
larger effects than in competitive markets. See MS-LM-008 (minimum wage in granular markets).
\end{axiom}

\subsection{10C Mapping for Monopsony}

\begin{table}[h]
\centering
\begin{tabular}{lll}
\textbf{10C Element} & \textbf{Monopsony Contribution} & \textbf{Parameter} \\
\hline
CORE-WHO (AAA) & Workers in concentrated markets & $L_j$ employment shares \\
CORE-WHAT (C) & Wage utility below MRP & $\mu = w/MRP$ \\
CORE-HOW (B) & Granularity $\times$ strategic interaction & $\gamma > 0$ \\
CORE-WHEN (V) & $\Psi_I$: regulatory context & HHI, noncompete laws \\
CORE-WHERE (BBB) & PAR-LM-006 through PAR-LM-009 & Structural estimates \\
CORE-HIERARCHY (HI) & Firm-level $\to$ market-level power & Granularity channel \\
\end{tabular}
\caption{Berger, Herkenhoff \& Mongey (2026) mapped to 10C framework}
\end{table}

\textbf{Theory Support:} MS-LM-006 (Granular Monopsony), MS-LM-007 (Oligopsonistic Wage Setting),
MS-LM-008 (Minimum Wage in Granular Markets). Case: CAS-955.

\subsection{Worked Example: Monopsony Impact Assessment}

Consider a US commuting zone with 5 major employers in manufacturing:

\begin{enumerate}
    \item \textbf{Measure concentration:} $HHI_L = \sum s_j^2 = 3{,}400$ (above 2,500 threshold)
    \item \textbf{Estimate markdown:} $\mu \approx 0.72$ (PAR-LM-006, manufacturing context)
    \item \textbf{Calculate wage gap:} Worker earning \$52,000 has $MRP \approx \$72,200$
    \item \textbf{Welfare loss:} 7.6\% of lifetime consumption (PAR-LM-008)
    \item \textbf{Policy intervention:}
    \begin{itemize}
        \item Ban noncompetes $\to$ +4\% wages $\to$ \$54,080 (PAR-LM-009)
        \item Antitrust enforcement $\to$ reduce $HHI$ below 2,500 $\to$ wage gains up to 30\%
    \end{itemize}
    \item \textbf{$\Psi$-Context:} Effects larger where $\Psi_I$ = ``weak labor regulation'' and $\Psi_S$ = ``few alternative employers''
\end{enumerate}


% ============================================================================
% OCCUPATIONAL LICENSING (Level 5 Integration: johnson2026occupational)
% ============================================================================
\section{Occupational Licensing in the United States}
\label{sec:dc-licensing}

\subsection{Johnson (2026): Institutional Barriers and Rent-Seeking}

\textbf{Citation.} Johnson, Janna E. (2026).
``Occupational Licensing in the United States.''
\textit{Journal of Economic Perspectives}, 40(1), 167--190. \cite{johnson2026occupational}

\textbf{Full text:} \texttt{data/paper-texts/PAP-johnson2026occupational.md} (L3).

\begin{tcolorbox}[colback=blue!5!white, colframe=blue!75!black, title=Occupational Licensing: Key Findings]
The paper documents licensing as the \textbf{dominant institutional barrier} ($\Psi_I$) in US labor markets:
\begin{enumerate}
    \item \textbf{Growth:} 24\% of the US workforce now requires an occupational license, up from just 5\% in 1950---a fivefold expansion in 74 years (PAR-LM-014).
    \item \textbf{Wage Premium:} Licensed workers earn approximately 15\% more than comparable unlicensed workers (PAR-LM-015).
    \item \textbf{Late Adopters:} Analysis of states that adopted licensing later shows a causal +4\% wage effect 25 years after adoption.
    \item \textbf{Professional Organizations:} Founding a professional organization increases the probability of licensing by 20 percentage points---suggesting \textbf{incumbent rent-seeking} ($\Psi_S$) as the primary driver, not consumer protection.
    \item \textbf{Composition:} Teachers and nurses account for 20\% of all licensed workers.
    \item \textbf{Quality:} Evidence that licensing improves service quality is mixed/weak; evidence of entry restriction and reduced geographic mobility is strong.
    \item \textbf{Scope:} 800+ occupations licensed with substantial cross-state variation in $\Psi_I$.
\end{enumerate}
\end{tcolorbox}

\begin{axiom}[UNMAPPED\_DC-14: Licensing as Institutional Rent Extraction]
\label{ax:dc14}
Occupational licensing is primarily driven by incumbent rent-seeking ($\Psi_S$), not consumer protection:
\begin{equation}
P(\text{licensing}_t) = \alpha + 0.20 \cdot \mathbb{1}[\text{professional\_org\_founded}_{t-k}] + \mathbf{X}\beta + \varepsilon
\end{equation}
The founding of a professional organization increases licensing probability by 20pp, demonstrating that $\Psi_S$ (social/organizational context) systematically \textbf{creates} $\Psi_I$ (institutional barriers) for rent extraction, disguised as consumer protection.

\textbf{EBF Connection:} CORE-HOW: This is a case where social actors ($\Psi_S$: professional organizations) systematically reshape institutional context ($\Psi_I$: licensing requirements). The complementarity between $\Psi_S$ and $\Psi_I$ is positive and causal. Theory: MS-LM-011. Case: CAS-958.
\end{axiom}

\begin{axiom}[UNMAPPED\_DC-15: Licensing and Geographic Mobility]
\label{ax:dc15}
State-specific licensing creates a mobility friction ($\Psi_I$) distinct from noncompetes ($\Psi_C$):
\begin{itemize}
    \item Noncompetes (Starr): Operate through \textbf{behavioral} channels---perception matters more than law
    \item Licensing (Johnson): Operates through \textbf{institutional} channels---legal requirements are explicit
    \item Both reduce mobility but through different $\Psi$ mechanisms
\end{itemize}
Together with employer concentration (Berger) and noncompetes (Starr), licensing completes the \textbf{three pillars of labor market power}: concentration-based, behavioral, and institutional.

\textbf{EBF Connection:} CORE-WHEN ($\Psi_I$): Workers must re-qualify when moving across state lines. This complements noncompete effects (CAS-957) and concentration effects (CAS-955/956) as a third mechanism reducing worker mobility and bargaining power.
\end{axiom}

\subsection{10C Mapping for Occupational Licensing}

\begin{table}[h]
\centering
\begin{tabular}{lll}
\textbf{10C Element} & \textbf{Johnson Contribution} & \textbf{Parameter/Reference} \\
\hline
CORE-WHAT (C) & Wage premium, entry costs & PAR-LM-015 (15\% premium) \\
CORE-WHERE (BBB) & Prevalence, premium, org effect & PAR-LM-014, PAR-LM-015 \\
CORE-WHEN (V) & $\Psi_I$ evolution 1950--2024 & 5\% $\to$ 24\% \\
CORE-WHO (AAA) & Teachers + nurses = 20\% licensed & Composition data \\
CORE-HOW (B) & $\Psi_S$ creates $\Psi_I$ (rent-seeking) & 20pp org effect \\
CORE-HIERARCHY (HI) & Multi-level: state, org, federal & 800+ occupations \\
\end{tabular}
\caption{Johnson (2026) mapped to 10C framework}
\end{table}

\textbf{Theory Support:} MS-LM-011 (Occupational Licensing as Labor Market Entry Barrier).
Case: CAS-958. Linked cases: CAS-955 (berger), CAS-956 (prager), CAS-957 (starr).

% ============================================================================
% SYNTHESIS: JEP 2026 LABOR MARKET POWER SYMPOSIUM
% ============================================================================
\section{Synthesis: Three Pillars of Labor Market Power}
\label{sec:dc-synthesis-jep2026}

The four JEP Winter 2026 papers form a coherent framework for understanding labor market power through complementary $\Psi$-mechanisms:

\begin{tcolorbox}[colback=green!5!white, colframe=green!75!black, title=JEP 2026 Symposium: Unified Framework]
\begin{tabular}{llll}
\textbf{Pillar} & \textbf{Paper} & \textbf{Mechanism} & \textbf{$\Psi$-Channel} \\
\hline
Concentration & Berger et al. & Granular monopsony & $\Psi_S$ (market structure) \\
Behavioral & Starr & Noncompete clauses & $\Psi_C$ (cognitive/perception) \\
Institutional & Johnson & Occupational licensing & $\Psi_I$ (legal requirements) \\
Enforcement & Prager & Antitrust in labor markets & $\Psi_I$ (regime change) \\
\end{tabular}

\textbf{Key Insight:} These mechanisms are \textbf{complementary} ($\gamma > 0$). In concentrated markets ($\Psi_S$) with noncompetes ($\Psi_C$) and licensing ($\Psi_I$), the combined effect on wages exceeds the sum of individual effects. Eliminating any single pillar captures less than one-third of the potential welfare gain.

\textbf{Parameters:} PAR-LM-006--015 (10 parameters). \textbf{Theories:} MS-LM-006--011 (6 theories). \textbf{Cases:} CAS-955--958 (4 cases).
\end{tcolorbox}


%%%%%%%%%%%%%%%%%%%%%%%%%%%%%%%%%%%%%%%%%%%%%%%%%%%%%%%%%%%%%%%%%%%%%%%%%%%%%%%
\subsubsection{2026: {From Asia, with Skills}}
%%%%%%%%%%%%%%%%%%%%%%%%%%%%%%%%%%%%%%%%%%%%%%%%%%%%%%%%%%%%%%%%%%%%%%%%%%%%%%%

\paragraph{Citation.}
Khanna, Gaurav (2026). ``{From Asia, with Skills}.'' \textit{Journal of Economic Perspectives}, 40, 215--240.

\paragraph{Core Finding.}
JEP Winter 2026 symposium paper. Full text available in data/paper-texts/PAP-khanna2026asia.md.

\paragraph{Key Parameters.}
\texttt{high_skill_share_india=0.78, high_skill_share_china=0.63, sector_software=0.38, visa_H1B_elasticity}

\paragraph{10C Integration.}
\begin{itemize}[nosep]
  \item \textbf{Domain}: [To be specified]
  \item \textbf{use\_for}: LIT-O, DOMAIN-LABOR, DOMAIN-MIGRATION, CORE-WHO, CORE-HOW, LIT-CARD
\end{itemize}


%%%%%%%%%%%%%%%%%%%%%%%%%%%%%%%%%%%%%%%%%%%%%%%%%%%%%%%%%%%%%%%%%%%%%%%%%%%%%%%
\subsubsection{2026: {Asian Immigration to the United States in Historical Perspe...}
%%%%%%%%%%%%%%%%%%%%%%%%%%%%%%%%%%%%%%%%%%%%%%%%%%%%%%%%%%%%%%%%%%%%%%%%%%%%%%%

\paragraph{Citation.}
Postel, Hannah M. (2026). ``{Asian Immigration to the United States in Historical Perspective}.'' \textit{Journal of Economic Perspectives}, 40, 191--214.

\paragraph{Core Finding.}
JEP Winter 2026 symposium paper. Full text available in data/paper-texts/PAP-postel2026asian.md.

\paragraph{Key Parameters.}
\texttt{population_growth_1960_2019=2700pct, policy_leverage_1965_act}

\paragraph{10C Integration.}
\begin{itemize}[nosep]
  \item \textbf{Domain}: [To be specified]
  \item \textbf{use\_for}: LIT-O, DOMAIN-MIGRATION, CORE-WHO, CORE-WHEN, LIT-CARD
\end{itemize}


% ============================================================================
% ANTITRUST ENFORCEMENT (Level 5 Integration: prager2026antitrust)
% ============================================================================
\section{Antitrust Enforcement in Labor Markets}
\label{sec:dc-antitrust}

\subsection{Prager (2026): From Product Markets to Labor Markets}

\textbf{Citation.} Prager, Elena (2026).
``Antitrust Enforcement in Labor Markets.''
\textit{Journal of Economic Perspectives}, 40(1), 115--138. \cite{prager2026antitrust}

\textbf{Full text:} \texttt{data/paper-texts/PAP-prager2026antitrust.md} (L3).

\begin{tcolorbox}[colback=blue!5!white, colframe=blue!75!black, title=Antitrust Enforcement: Key Findings]
The paper documents the dramatic expansion of antitrust enforcement into labor markets:
\begin{enumerate}
    \item \textbf{Enforcement Growth:} FTC/DOJ labor market enforcement actions grew 700\% in one decade---from 1 action (2010--14) to 7 actions (2020--24) (PAR-LM-010).
    \item \textbf{2023 Merger Guidelines:} First mention of labor market effects in their 55-year history---a fundamental $\Psi_I$ regime change.
    \item \textbf{US v.\ Bertelsmann (2022):} First merger blocked on labor market grounds (publishing industry), establishing enforcement precedent.
    \item \textbf{Cournot Wage Markdown:} Under Cournot competition, wage markdowns are proportional to HHI: $markdown_i = s_i / \varepsilon$ (PAR-LM-011).
    \item \textbf{Key Distinction:} Concentration-based labor market power (actionable via antitrust) vs.\ search friction-based power (not directly actionable).
\end{enumerate}
\end{tcolorbox}

\begin{axiom}[UNMAPPED\_DC-10: Cournot Wage Markdown under Antitrust]
\label{ax:dc10}
Under Cournot (quantity) competition in labor markets, firm $i$'s optimal wage markdown is:
\begin{equation}
\frac{MRP_i - w_i}{w_i} = \frac{s_i}{\varepsilon_S}
\end{equation}
where $s_i$ is firm $i$'s employment share and $\varepsilon_S$ is the labor supply elasticity. This directly links HHI to wage suppression: higher concentration $\Rightarrow$ larger markdowns.

\textbf{EBF Connection:} Extends UNMAPPED\_DC-7 (granular monopsony) with an antitrust-actionable mechanism. Links to CORE-HOW: the functional form connecting market structure to wage outcomes is proportional, not threshold-based. Theory: MS-LM-009. Parameter: PAR-LM-011.
\end{axiom}

\begin{axiom}[UNMAPPED\_DC-11: Institutional Evolution of Labor Antitrust]
\label{ax:dc11}
The enforcement regime ($\Psi_I$) for labor market antitrust evolved in three phases:
\begin{enumerate}
    \item \textbf{Pre-2016:} No-poach agreements treated as employment law, not antitrust
    \item \textbf{2016--2022:} DOJ begins criminal prosecution of no-poach/wage-fixing
    \item \textbf{2023+:} Merger Guidelines explicitly incorporate labor market analysis
\end{enumerate}
\textbf{EBF Connection:} CORE-WHEN ($\Psi_I$ regime change): the institutional context shifted from ``labor markets are not antitrust territory'' to ``labor market effects are a first-class concern in merger review.'' Case: CAS-956.
\end{axiom}

\subsection{10C Mapping for Antitrust Enforcement}

\begin{table}[h]
\centering
\begin{tabular}{lll}
\textbf{10C Element} & \textbf{Prager Contribution} & \textbf{Parameter/Reference} \\
\hline
CORE-HOW (B) & Cournot wage markdown $\propto$ HHI & PAR-LM-011 \\
CORE-WHERE (BBB) & Enforcement data, markdown formula & PAR-LM-010, PAR-LM-011 \\
CORE-WHEN (V) & $\Psi_I$ regime change (2023 Guidelines) & CAS-956 \\
CORE-HIERARCHY (HI) & Conduct $\to$ Merger $\to$ Policy levels & MS-LM-009 \\
\end{tabular}
\caption{Prager (2026) mapped to 10C framework}
\end{table}

\textbf{Theory Support:} MS-LM-009 (Antitrust Enforcement in Labor Markets).
Case: CAS-956. Linked cases: CAS-955 (berger), CAS-957 (starr), CAS-958 (johnson).


% ============================================================================
% NONCOMPETE CLAUSES (Level 5 Integration: starr2026economics)
% ============================================================================
\section{The Economics of Noncompete Clauses}
\label{sec:dc-noncompete}

\subsection{Starr (2026): Behavioral Mechanisms of Labor Market Restraints}

\textbf{Citation.} Starr, Evan (2026).
``The Economics of Noncompete Clauses.''
\textit{Journal of Economic Perspectives}, 40(1), 139--166. \cite{starr2026economics}

\textbf{Full text:} \texttt{data/paper-texts/PAP-starr2026economics.md} (L3).

\begin{tcolorbox}[colback=blue!5!white, colframe=blue!75!black, title=Noncompete Clauses: Key Findings]
The paper reveals that noncompetes operate through \textbf{behavioral} channels as much as legal ones:
\begin{enumerate}
    \item \textbf{Prevalence:} 18\% of US workers are bound by noncompetes; 30\% of low-wage firms require them (PAR-LM-012).
    \item \textbf{Ban Wage Effect:} State-level bans (Oregon, Hawaii, Washington) increase wages 3--4\% via natural experiments (PAR-LM-013).
    \item \textbf{RCT Evidence:} An experiment with 14,000 job offers found noncompetes reduce the probability of working for a competitor by 57\%.
    \item \textbf{Contract Non-Reading:} One-third of workers skip reading the noncompete clause buried on page 7 of the contract ($\Psi_C$: cognitive overload at hiring).
    \item \textbf{Enforceability Paradox:} Noncompete wage premia are \textbf{larger} in states where they are unenforceable---the \textit{in terrorem} effect means $\Psi_C$ (perceived constraint) dominates $\Psi_I$ (legal enforceability).
    \item \textbf{California:} Ban since 1872 is credited for Silicon Valley's innovation ecosystem---a long-run $\Psi_I$ effect on knowledge spillovers.
    \item \textbf{Compensation:} Requiring a noncompete reduces total compensation by 12\%; signing one reduces it by 4.2\%.
\end{enumerate}
\end{tcolorbox}

\begin{axiom}[UNMAPPED\_DC-12: In Terrorem Effect (Behavioral $>$ Legal)]
\label{ax:dc12}
The behavioral effect of noncompetes exceeds the legal effect:
\begin{equation}
\text{Mobility Reduction} = f(\Psi_C^{perceived}) > f(\Psi_I^{actual})
\end{equation}
Workers in \textbf{non-enforcing} states with noncompetes are less mobile than predicted by legal enforceability alone. The perception of constraint ($\Psi_C$) is a stronger behavioral mechanism than the actual institutional constraint ($\Psi_I$).

\textbf{EBF Connection:} CORE-AWARE: Workers have \textbf{absent} or \textbf{biased} awareness of their legal rights. The one-third who don't read the clause operate in \textit{pre-contemplation} (CORE-STAGE). Policy implication: banning noncompetes works through \textbf{both} legal AND behavioral channels simultaneously---removing $\Psi_I$ also eliminates the $\Psi_C$ anchor. Theory: MS-LM-010. Parameters: PAR-LM-012, PAR-LM-013.
\end{axiom}

\begin{axiom}[UNMAPPED\_DC-13: Cognitive Overload at Contract Signing]
\label{ax:dc13}
At the moment of hiring, workers face a \textbf{cognitive overload} configuration:
\begin{itemize}
    \item Bargaining power is weakest (offer already extended)
    \item Cognitive load is highest (new job, relocation, benefits)
    \item Noncompete is buried on page 7 of the contract
    \item 33\% of workers don't read it at all
\end{itemize}
\begin{equation}
P(\text{sign without reading}) \approx 0.33 \quad \text{when} \quad \Psi_C = \text{``cognitive\_overload\_hiring''}
\end{equation}
\textbf{EBF Connection:} CORE-STAGE: Signing occurs at \textit{pre-contemplation/action} with no evaluation phase. CORE-READY: Willingness to negotiate is near zero due to asymmetric power. Case: CAS-957.
\end{axiom}

\subsection{10C Mapping for Noncompete Clauses}

\begin{table}[h]
\centering
\begin{tabular}{lll}
\textbf{10C Element} & \textbf{Starr Contribution} & \textbf{Parameter/Reference} \\
\hline
CORE-HOW (B) & Noncompetes as wage suppression mechanism & PAR-LM-012, PAR-LM-013 \\
CORE-WHAT (C) & Compensation effects ($-12\%$, $-4.2\%$) & Survey regression \\
CORE-AWARE (AU) & In terrorem effect; contract non-reading & $\Psi_C > \Psi_I$ \\
CORE-READY (AV) & Bargaining power at hiring $\approx 0$ & CAS-957 \\
CORE-STAGE (AW) & Signing at pre-contemplation & No evaluation phase \\
CORE-WHEN (V) & California ban since 1872 & Long-run $\Psi_I$ \\
\end{tabular}
\caption{Starr (2026) mapped to 10C framework}
\end{table}

\textbf{Theory Support:} MS-LM-010 (Noncompete Clauses as Labor Market Restraints).
Case: CAS-957. Linked cases: CAS-955 (berger), CAS-956 (prager), CAS-958 (johnson).

\subsection{Worked Example: Noncompete Policy Analysis}

Consider a state legislature evaluating a noncompete ban for workers earning below \$100,000:

\begin{enumerate}
    \item \textbf{Baseline:} 18\% of workers bound (PAR-LM-012), including 30\% of low-wage firms
    \item \textbf{Predicted wage effect:} +3--4\% for affected workers (PAR-LM-013, from Oregon/Hawaii evidence)
    \item \textbf{Behavioral channel:} In terrorem effect eliminated---workers \textbf{know} they can leave
    \item \textbf{Mobility effect:} 57\% reduction in competitor employment eliminated (RCT evidence)
    \item \textbf{Innovation:} Long-run knowledge spillover gains (California precedent)
    \item \textbf{$\Psi$-Context:} Effects larger where $\Psi_C$ = ``low legal literacy'' and $\Psi_S$ = ``employer-dominated norms''
\end{enumerate}

% === END AUTO-GENERATED ENTRIES ===
\section{Summary}

\begin{tcolorbox}[colback=gray!10!white, colframe=gray!75!black, title=Key Takeaways]
\begin{enumerate}
    \item \textbf{Natural Experiments:} Exploit policy variation for causal identification
    \item \textbf{Minimum Wage:} Near-zero employment effects for moderate increases
    \item \textbf{Immigration:} Labor markets absorb large inflows with small wage effects
    \item \textbf{Returns to Education:} 8-12\% per year, robust to endogeneity corrections
    \item \textbf{Firm Effects:} 20-25\% of wage variance, increasing sorting over time
    \item \textbf{ALMP:} Training works long-run; public employment doesn't work
\end{enumerate}
\end{tcolorbox}

% ============================================================================
% GLOSSARY
% ============================================================================
\section{Glossary}

See Appendix G (Glossary) for standard EBF terms. Card-specific terms:

\begin{description}
    \item[Design-Based Research] Causal inference prioritizing identification through research design
    \item[Difference-in-Differences] Method comparing treatment-control differences before/after intervention
    \item[Firm Fixed Effect] Time-invariant wage premium associated with a specific employer
    \item[Monopsony] Labor market with wage-setting power by employers
    \item[Natural Experiment] Naturally occurring variation that mimics random assignment
\end{description}

% ============================================================================
% CRITICAL FOUNDATIONS
% ============================================================================
\section{Critical Foundations}

\begin{enumerate}
    \item Card, D. \& Krueger, A. (1994). Minimum Wages and Employment. \textit{AER}. [Revolutionary]
    \item Card, D. (1990). The Impact of the Mariel Boatlift. \textit{ILRR}. [Natural Experiment Classic]
    \item Card, D. (1999). The Causal Effect of Education on Earnings. \textit{Handbook}. [Survey]
    \item Card, D. et al. (2013). Workplace Heterogeneity and Wage Inequality. \textit{QJE}. [Firm Effects]
    \item Card, D. et al. (2018). What Works? ALMP Meta-Analysis. \textit{JEEA}. [Policy]
\end{enumerate}

% ============================================================================
% OPEN ISSUES
% ============================================================================
\section{Open Issues}

\begin{enumerate}
    \item \textbf{External validity:} Do natural experiment results generalize to other contexts?
    \item \textbf{Minimum wage thresholds:} At what level do negative employment effects emerge?
    \item \textbf{Immigration long-run:} What are the dynamic effects beyond short-run adjustment?
    \item \textbf{Firm effects mechanisms:} Why do some firms pay more? Rent-sharing vs. sorting?
\end{enumerate}

% ============================================================================
% REFERENCES
% ============================================================================
\section{References}

\nocite{bcm_master}

For complete bibliography, see Master Bibliography (\texttt{bibliography/master\_references.bib}).

Key Card references integrated in this appendix:
\begin{itemize}[nosep]
    \item card1994minimum, card1990mariel, card1995myth
    \item card1999causal, card2001estimating, card2012inequality
    \item card2012peer, card2015active, card2018works
    \item card2022nobel
\end{itemize}

\end{document}
